\documentclass[11pt]{amsart}
\usepackage{amsmath,amssymb,amsthm}
\usepackage[margin=1.1in]{geometry}
\usepackage{booktabs}
\usepackage[colorlinks=true,allcolors=blue]{hyperref}
\usepackage[capitalise,nameinlink]{cleveref}

\newtheorem{theorem}{Theorem}[section]
\newtheorem{proposition}[theorem]{Proposition}
\newtheorem{lemma}[theorem]{Lemma}
\newtheorem{corollary}[theorem]{Corollary}
\theoremstyle{definition}
\newtheorem{definition}[theorem]{Definition}
\theoremstyle{remark}
\newtheorem{remark}[theorem]{Remark}

\newcommand{\E}{\mathbb{E}}
\newcommand{\R}{\mathbb{R}}
\newcommand{\Simp}{\Delta_4}
\newcommand{\one}{\mathbf{1}}

\title{The Kesten--Stigum bound is tight for the four-state Potts model on the binary tree}
\author{Yitzchak Shmalo}
\date{\today}

\usepackage{xurl}
\begin{document}

\begin{abstract}
Consider the broadcast of a uniformly random colour from the root of the infinite binary tree, where
each edge copies its parent's colour with probability $\lambda+(1-\lambda)/4$ and otherwise chooses one
of the other three colours uniformly. We prove that for $0<\lambda\le 1/\sqrt2$ the colours at depth
$n$ carry asymptotically no information about the root colour. Since reconstruction is possible for
$2\lambda^2>1$ by the theorem of Kesten and Stigum, the Kesten--Stigum bound is the reconstruction
threshold for the ferromagnetic four-state Potts channel on the binary tree. A preprint of OpenAI
establishes this threshold on every regular tree; the proof given here is different. It is computer
assisted. A certified upper bound on the second moment of the root posterior, obtained from an exact
finite-state upgrade of the belief-propagation recursion, places the limiting posterior law in a small
neighbourhood of the uniform law; an exact polynomial certificate then shows that no nontrivial
belief-propagation fixed point lies in that neighbourhood.
\end{abstract}

\maketitle

\section{Introduction}

Let $T$ be the infinite rooted binary tree. Fix $q\ge2$ and $\lambda\in(0,1)$ and let
$M_\lambda=\lambda I+(1-\lambda)J$ be the $q\times q$ Potts channel, where $J$ has every entry $1/q$.
The broadcast process assigns the root $\rho$ a uniform colour $\sigma_\rho\in[q]$ and passes colours
down each edge independently through $M_\lambda$. Write $\sigma_{L_n}$ for the colours at depth $n$.
The process is \emph{reconstructible} if
\[
\lim_{n\to\infty}\E\,\big\|\,\mathbb P(\sigma_\rho=\cdot\mid\sigma_{L_n})-q^{-1}\one\,\big\|_{\mathrm{TV}}>0,
\]
and nonreconstructible otherwise. Kesten and Stigum~\cite{KS66} showed that reconstruction holds when
$d\lambda^2>1$ on the $d$-ary tree, and for $q=2$ this is sharp~\cite{BRZ95,EKPS00}. For larger $q$ the
bound can fail to be sharp: Mossel~\cite{Mos01} showed that for $q$ large reconstruction holds below it,
and Sly~\cite{Sly11} proved that the bound is not tight for $q\ge5$ and every large $d$, and that it is
tight for $q=3$ and large $d$; for $q=3$ tightness holds on the $d$-ary tree for every $d\ge2$~\cite{Shm3}.
Mossel, Sly and Sohn~\cite{MSS22} extended tightness to $q=4$ for large $d$.
Ricci-Tersenghi, Semerjian and Zdeborov\'a~\cite{RSZ19} predicted, from an expansion of the cavity
equations, that for the ferromagnetic four-state Potts model the Kesten--Stigum bound is tight for every
degree distribution. Bhatnagar and Maneva~\cite{BM11} developed a computational method for small trees
and recorded the gap that it leaves. The preprint~\cite{OAI4} establishes the ferromagnetic four-state
threshold on every regular tree and on every Poisson Galton--Watson tree: a polynomial constraint on the
law of the posterior is preserved by the recursion, and it yields an inequality
$m_{n+1}\le d\lambda^2m_n-cm_n^3$ for the expected quadratic information $m_n$ at depth $n$.

The main result of this paper is the binary-tree case of that result, proved by a different route.

\begin{theorem}\label{thm:main}
For the four-state Potts channel on the binary tree, the broadcast process is nonreconstructible for
every $0<\lambda\le1/\sqrt2$. In particular the reconstruction threshold equals the Kesten--Stigum bound
$\lambda=1/\sqrt2$.
\end{theorem}

The proof has two parts of different kinds. The first is an upper bound on the second moment
$x_n=\E\sum_i(p_i-1/4)^2$ of the root posterior $p=\mathbb P(\sigma_\rho=\cdot\mid\sigma_{L_n})$, valid
for $n\ge500$ (\cref{prop:grid}). It is obtained by running the belief-propagation recursion on an
information upgrade of the posterior law that lives on a finite grid, in exact integer arithmetic. The
second is a fixed-point exclusion. The limit of the posterior laws is a fixed point of the
belief-propagation map (\cref{lem:fixed}), and we construct an explicit function $K(a,b;m)$ of two
messages and a scalar whose expectation at any fixed point equals $256(x-x_{\mathrm{out}})$ minus a
positive multiple of $m^3$ (\cref{prop:expect}). A finite family of exact polynomial certificates shows
$K\ge0$ on the region permitted by the first part (\cref{prop:kernel}). At a fixed point
$x=x_{\mathrm{out}}$, so the second moment is zero.

Neither part alone suffices. The grid bound $x_n\le0.0453$ is far from zero, and an exclusion argument
based on a local expansion around the uniform law covers a much smaller neighbourhood. The kernel $K$ is chosen
so that its region of validity reaches the grid bound. It combines the drift identity for the second
moment with a family of $31$ functions whose expectations vanish at every colour-symmetric law, and a
sign condition on the third moment that holds at every fixed point (\cref{prop:cubic}).

\section{The recursion}\label{sec:recursion}

Throughout $q=4$. Let $\Simp=\{p\in\R^4:p_i\ge0,\ \sum_ip_i=1\}$. A posterior $p$ is encoded by the
\emph{message} $a=\lambda(4p-\one)$, which lies in $\lambda(4\Simp-\one)$ and has $\sum_ia_i=0$. If the
two children of a vertex send messages $a$ and $b$, the posterior of the vertex is
\[
P_i=\frac{(1+a_i)(1+b_i)}{\sum_j(1+a_j)(1+b_j)},
\qquad
4P-\one=\frac{H}{s},
\]
where
\[
c=a\cdot b,\qquad s=1+\tfrac c4,\qquad H_i=a_i+b_i+a_ib_i-\tfrac c4 .
\]
The outgoing message is $\lambda H/s$. Let $\mu$ be the law of the message at a vertex, taken
unconditionally on the root colour. The law $\mu$ is invariant under colour permutations and has mean
zero. The two children's observations are independent given the parent's colour but not
unconditionally; their joint law has density $s(a,b)$ with respect to $\mu\otimes\mu$. We call
$\mu\otimes\mu$ the base measure and write $\E$ for expectations with respect to it, so the
belief-propagation map $\mathcal T_\lambda$ is
\begin{equation}\label{eq:bp}
\int f\,d(\mathcal T_\lambda\mu)=\E\Big[s(a,b)\,f\big(\lambda H(a,b)/s(a,b)\big)\Big].
\end{equation}
We use the notation
\[
A=|a|^2,\quad B=|b|^2,\quad t_a=\textstyle\sum_ia_i^3,\quad t_b=\sum_ib_i^3,\quad
p_{ij}=\sum_k a_k^ib_k^j,\quad \ell=p_{21}+p_{12},
\]
and the moments of $\mu$,
\[
m=\E A,\quad T=\E t_a,\quad V=\E A^2,\quad Q=\E\textstyle\sum_ia_i^4 .
\]
The second moment of the posterior is $x=\E\sum_i(p_i-1/4)^2=m/(16\lambda^2)$; at $\lambda=1/\sqrt2$,
$m=8x$. From \eqref{eq:bp},
\begin{equation}\label{eq:moments-out}
m_{\mathrm{out}}=\lambda^2\,\E\Big[\frac Ns\Big],\quad
T_{\mathrm{out}}=\lambda^3\,\E\Big[\frac J{s^2}\Big],\quad
V_{\mathrm{out}}=\lambda^4\,\E\Big[\frac{N^2}{s^3}\Big],
\qquad N=\textstyle\sum_iH_i^2,\ J=\sum_iH_i^3 .
\end{equation}

\section{The canonical bound}\label{sec:grid}

Let $\mu_n$ be the law of the root message given depth $n$. Its second moment is $16\lambda^2x_n$.

\begin{proposition}\label{prop:grid}
Let
\[
U=\frac{1375940705703331}{30399297484750848}<0.04527 .
\]
For every $0<\lambda\le707107/10^6$ and every $n\ge500$, $x_n\le U$.
\end{proposition}

The proof compares experiments in the sense of Blackwell. A law of posteriors with uniform mean is an
experiment about the root colour, and replacing it by a more informative experiment can only increase
$x$, because $x$ is the expectation of a convex function of the posterior. Belief propagation preserves
the order: if each child's experiment is replaced by a more informative one, so is the parent's. A
posterior $p$ may be replaced by a random posterior with barycentre $p$ (reveal additional
information); mass may be moved to the point posterior of a colour chosen by a colour-equivariant
rule (full revelation). The finite computation starts at full revelation, which is the depth-zero
law, and alternates one exact belief-propagation step with two such upgrades.

\begin{lemma}[symmetric lattice splitting]\label{lem:split}
Let $N\ge1$ and $p\in\Simp$. Write $Np_i=f_i+r_i$ with $f_i\in\mathbb Z$ and $0\le r_i<1$, and let
$R=\sum_ir_i\in\{0,1,2,3\}$. Define a law on $N^{-1}\mathbb Z^4\cap\Simp$ as follows. If $R=0$, the
point mass at $f/N$. If $R=1$, the point $(f+e_i)/N$ with weight $r_i$. If $R=3$, the point
$(f+\one-e_i)/N$ with weight $1-r_i$. If $R=2$, with $h=\min_i\min(r_i,1-r_i)$, the point
$(f+e_i+e_j)/N$ with weight
\[
w_{ij}=\max(r_i+r_j-1,0)+h/3\qquad(i<j).
\]
In every case the weights are nonnegative, sum to one, have barycentre $p$, and depend equivariantly
on $p$ under colour permutations.
\end{lemma}

\begin{proof}
The cases $R\in\{0,1,3\}$ are immediate. Let $R=2$ and sort the residuals $r_{(1)}\le r_{(2)}\le
r_{(3)}\le r_{(4)}$. Then $r_{(1)}+r_{(2)}\le1$, $r_{(1)}+r_{(3)}\le1$, $r_{(2)}+r_{(4)}\ge1$ and
$r_{(3)}+r_{(4)}\ge1$. If $r_{(1)}+r_{(4)}\ge1$, the positive excesses $r_i+r_j-1$ are those of the
three pairs containing the largest residual, and for each coordinate $i$ the excesses of the pairs
containing $i$ add up to $r_i-(1-r_{(4)})$. If $r_{(1)}+r_{(4)}\le1$, the positive excesses are those
of the three pairs not containing the smallest residual, and the incident excesses add up to
$r_i-r_{(1)}$. In both cases they add up to $r_i-h$, so adding $h/3$ to each of the three pairs at $i$
gives total weight $r_i$ at coordinate $i$. Summing over $i$ counts each pair twice, so the weights add
up to $R/2=1$. The barycentre is $(f+\sum_{i<j}w_{ij}(e_i+e_j))/N=(f+r)/N=p$.
\end{proof}

\begin{proof}[Proof of \cref{prop:grid}]
Write $\lambda^*=707107/10^6$; note $2(\lambda^*)^2>1$. Take $N=36$ and work with laws supported on
$\{v/36:v\in\mathbb Z_{\ge0}^4,\ \sum v_i=36\}$, grouped into the $478$ orbits of the colour
permutations. Starting from the point posterior, one step consists of the exact two-child update
\eqref{eq:bp} at $\lambda^*$ applied to the orbit law, the splitting of \cref{lem:split} applied to
every resulting posterior, and grouping of equal posteriors. The orbit masses are then rounded down to
integer multiples of $2^{-48}$ and the deficit is placed on the orbit of full revelation. All
transition probabilities have the common denominator $288(36\cdot10^6)^2$, and every quantity is an
exact integer.

By \cref{lem:split}, splitting is an upgrade. Rounding down and moving the deficit to full revelation
is an upgrade because each orbit is colour-symmetric with uniform mean. By induction the computed law
at step $n$ is more informative than $\mu_n$ at $\lambda^*$. After $500$ steps its second moment equals
$U$ exactly. Since $x_n$ is nonincreasing in $n$ ($\sigma_{L_{n+1}}$ is obtained from $\sigma_{L_n}$ by
independent randomness, so it is a garbling of $\sigma_{L_n}$), the bound holds for all $n\ge500$. For $\lambda<\lambda^*$ it follows from
\cref{lem:garble}.
\end{proof}

\begin{lemma}\label{lem:garble}
For $0<\lambda'<\lambda<1$ and every $n$, $x_n(\lambda')\le x_n(\lambda)$.
\end{lemma}

\begin{proof}
Put $\theta=\lambda'/\lambda$. On each edge draw independently a permutation $R_v$ of the colours,
equal to the identity with probability $\theta$ and uniform otherwise, and set $G_\rho=\mathrm{id}$,
$G_v=G_{\mathrm{parent}(v)}\circ R_v$. If $\sigma$ is the $\lambda$-broadcast, then $\sigma'_v=G_v(\sigma_v)$
is the $\lambda'$-broadcast, because averaging over $R_v$ gives the transition
$M_\lambda M_\theta=M_{\lambda'}$, which does not depend on $G_{\mathrm{parent}(v)}$. The root colour
is unchanged and $\sigma'_{L_n}$ is a function of $\sigma_{L_n}$ and independent randomness, so the
$\lambda'$ observation is a garbling of the $\lambda$ observation and Jensen's inequality gives the
claim.
\end{proof}

\section{The limit law}\label{sec:limit}

\begin{lemma}\label{lem:fixed}
Fix $\lambda\in(0,1)$. The laws $\mu_n$ converge weakly to a colour-symmetric law $\mu_\infty$ with
mean zero, $\mu_\infty=\mathcal T_\lambda\mu_\infty$, and $16\lambda^2x_n\to\int|a|^2\,d\mu_\infty$.
The process is nonreconstructible if and only if $\mu_\infty$ is the point mass at $0$.
\end{lemma}

\begin{proof}
Let $\mathcal F_n$ be generated by all colours at depth at least $n$. By the Markov property
$\mathbb P(\sigma_\rho=i\mid\mathcal F_n)=\mathbb P(\sigma_\rho=i\mid\sigma_{L_n})$, and the
$\sigma$-fields decrease, so L\'evy's downward theorem gives almost sure and $L^1$ convergence of the
posteriors. Hence $\mu_n$ converges weakly and the second moments converge. Symmetry and the mean pass
to the limit. In \eqref{eq:bp} the density $s$ is at least $1-\lambda^2>0$ and the integrand is a
bounded continuous function of $(a,b)$, so $\mathcal T_\lambda$ is weakly continuous, and
$\mu_{n+1}=\mathcal T_\lambda\mu_n$ gives the fixed-point equation. The last statement follows from
$\E\|p-\tfrac14\one\|_{\mathrm{TV}}\le\sqrt{x_n}$.
\end{proof}

\begin{lemma}[support]\label{lem:support}
Let $\varrho=677/1000$ and $0<\lambda\le1/\sqrt2$. Every fixed point $\mu$ of $\mathcal T_\lambda$ is
supported in $\varrho(4\Simp-\one)$.
\end{lemma}

\begin{proof}
A zero-sum vector lies in $r(4\Simp-\one)$ exactly when its coordinates are at least $-r$. Suppose
$a,b\in r(4\Simp-\one)$ with $0\le r<1$. For fixed $b$ the coordinate $P_i$ is a linear-fractional
function of $a$ with positive denominator, so its minimum over the simplex is attained at a vertex,
and likewise in $b$. At a pair of vertices the smallest coordinate of $P$ is $(1-r)^2/(4(1+3r^2))$ when
the two favoured colours agree and $(1-r)^2/(4(1-r^2))$ otherwise. Hence the outgoing message has
coordinates at least $-f(r)$ with $f(r)=2\ell r(1+r)/(1+3r^2)$, $\ell=707107/10^6$. The function $f$ is
increasing on $[0,1]$. Every message lies in $\ell(4\Simp-\one)$, and since $\mu=\mathcal T_\lambda^k\mu$
for every $k$, $\mu$ is supported in $r_k(4\Simp-\one)$ where $r_0=\ell$ and $r_{k+1}=f(r_k)$. The
sequence decreases to the fixed point $0.6757\ldots$ of $f$, and an exact rational computation shows
$r_k\le\varrho$ after finitely many steps.
\end{proof}

\section{The third moment at a fixed point}\label{sec:cubic}

\begin{proposition}\label{prop:cubic}
Let $0<\lambda\le1/\sqrt2$ and let $\mu$ be a colour-symmetric fixed point of $\mathcal T_\lambda$ with
second moment $m$ and third moment $T$. Then
\[
\tfrac{69763}{500000}\,m^2\le\big(\lambda^{-3}-2+\tfrac m6\big)\,T\le\tfrac76\,m^2 .
\]
In particular $T\ge0$.
\end{proposition}

\begin{proof}
For independent colour-symmetric $a,b$ supported in $\ell(4\Simp-\one)$ we show
\begin{equation}\label{eq:qcf}
\tfrac{69763}{500000}\,m_am_b\le\E\Big[\frac J{s^2}\Big]-\big(1-\tfrac{m_b}{12}\big)T_a-\big(1-\tfrac{m_a}{12}\big)T_b
\le\tfrac76\,m_am_b .
\end{equation}
With $a=b$ in law and $T_{\mathrm{out}}=T$, \eqref{eq:moments-out} turns \eqref{eq:qcf} into the
statement, and $\lambda^{-3}\ge2\sqrt2>2$ gives $T\ge0$.

For a coset $C$ of the Klein subgroup $\{e,(01)(23),(02)(13),(03)(12)\}$ in $S_4$ write $b_\pi$ for $b$
with coordinates permuted by $\pi$, and put
\[
G_C=\frac1{|C|}\sum_{\pi\in C}\frac{J(a,b_\pi)}{s(a,b_\pi)^2}-t_a-t_b+\frac{Bt_a+At_b}{12},\qquad
R_C=\frac3{|C|}\sum_{\pi\in C}(a\cdot b_\pi)^2-AB .
\]
For each of the six cosets we prove the pointwise inequality
\[
\beta\,AB\le G_C+R_C/2\le M_C\,AB,\qquad \beta=69763/500000,
\]
with constants $M_C$ whose average is $3498859/3000000<7/6$. Sort $a$ and parametrise both messages
by barycentric coordinates $z,w\ge0$ on the ordered simplex with vertices
\[
\ell(-1,-1,-1,3),\quad \ell(-1,-1,1,1),\quad \ell(-1,\tfrac13,\tfrac13,\tfrac13),\quad 0 .
\]
After clearing the positive
denominators $s_\pi\ge1-\ell^2$ and the homogenising factors, $G_C+R_C/2$ and $AB$ become polynomials of bidegree $(11,11)$ in $(z,w)$
with the same $129600$ monomials, and the two inequalities hold coefficient by coefficient. The six
cosets have the same size, so the average of $G_C+R_C/2$ over them lies between $\beta AB$ and
$\big(\tfrac16\sum_C M_C\big)AB\le\tfrac76AB$, and it is the average over all of $S_4$. Averaging
over cosets, using colour symmetry of each law separately, and taking expectations gives
\eqref{eq:qcf}, since the average of $R_C$ over all cosets is $3(a\cdot b)^2$ averaged over $S_4$ minus
$AB$, which is zero for zero-sum vectors.
\end{proof}

\section{The kernel}\label{sec:kernel}

\subsection{Moment identities}
Let $a,b$ be independent with the same colour-symmetric law of mean zero. For every scalar function
$w$ invariant under colour permutations,
\begin{equation}\label{eq:weighted}
\E[w(a)a_i]=0,\qquad \E[w(a)a_ia_j]=\tfrac13\E[w(a)A]\big(\delta_{ij}-\tfrac14\big),
\end{equation}
and the third-moment tensor is $\E a_i^3=T/4$, $\E a_i^2a_j=-T/12$, $\E a_ia_ja_k=T/12$ for distinct
indices. Conditioning on one of the two messages and applying these gives the expectations in
\cref{tab:moments}.

\begin{table}[h]
\centering
\begin{tabular}{ll@{\qquad}ll}
\toprule
expression & expectation & expression & expectation\\
\midrule
$c^2$ & $m^2/3$ & $p_{33}$ & $T^2/4$\\
$(A+B)c^2$ & $2mV/3$ & $p_{42}+p_{24}$ & $mQ/2$\\
$c\,\ell$ & $2mT/3$ & $c(p_{31}+p_{13})$ & $2mQ/3$\\
$c^3$ & $2T^2/3$ & $p_{21}^2+p_{12}^2$ & $\tfrac{2m}3(Q-V/4)$\\
$p_{22}$ & $m^2/4$ & $p_{21}p_{12}$ & $T^2/3$\\
$p_{32}+p_{23}$ & $mT/2$ & $(A+B)p_{22}$ & $mV/2$\\
$c\,p_{22}$ & $T^2/3$ & & \\
\bottomrule
\end{tabular}
\caption{Expectations under the base measure.}\label{tab:moments}
\end{table}

Expanding $N=A+B+2c+2\ell+p_{22}-c^2/4$ and $1/s=1-c/4+c^2/(16s)$, the table gives
$\E N=2m+m^2/6$ and $\E[Nc]=2m^2/3+4mT/3+T^2/6$. At $\lambda=1/\sqrt2$, \eqref{eq:moments-out}
then yields
\begin{equation}\label{eq:drift}
256\,(x-x_{\mathrm{out}})=\tfrac{16}3mT+\tfrac23T^2-\E\Big[\frac{Nc^2}{s}\Big].
\end{equation}

\subsection{Definition of the kernel}
For messages $a,b$ and a scalar $m\ge0$ let
\[
K(a,b;m)=\tfrac23t_at_b-\frac{Nc^2}{s}-\beta\,m\,AB+\sum_{j=0}^{32}g_j\,\phi_j(a,b;m),
\qquad \beta=10^{-6},
\]
where the functions $\phi_j$ are listed in \cref{tab:features} and the rational coefficients $g_j$ in
\cref{app:coeff}. Here $F_3=J/s^2$, $F_4=N^2/s^3$, $s_a=\sum_ia_i^4$, $s_b=\sum_ib_i^4$, $p_{31}$ stands
for $\sum_k(a_k^3b_k+a_kb_k^3)$, and $\kappa=c\,\ell-(Bt_a+At_b)/3$. The coefficients of $\phi_{25}$,
$\phi_{26}$, $\phi_{27}$ and $\phi_{28}$ are zero, so $K$ is affine in $m$.

\begin{table}[h]
\centering\small
\begin{tabular}{rl@{\qquad}rl}
\toprule
$j$ & $\phi_j$ & $j$ & $\phi_j$\\
\midrule
0 & $m(F_3-t_a-t_b)$ & 17 & $m\kappa$\\
1 & $AB(A+B)-m(F_4-A^2-B^2)$ & 18 & $m\big(p_{32}+p_{23}-(Bt_a+At_b)/4\big)$\\
2 & $(A+B)(c^2-AB/3)$ & 19 & $c\,p_{22}-t_at_b/3$\\
3 & $m(c^2-AB/3)$ & 20 & $p_{33}-t_at_b/4$\\
4 & $\kappa$ & 21 & $p_{42}+p_{24}-(Bs_a+As_b)/4$\\
5 & $c(A+B)$ & 22 & $mc(A^2+B^2)$\\
6 & $c(A^2+B^2)$ & 23 & $24c^4-\sum_{\pi\in S_4}(a\cdot b_\pi)^4$\\
7 & $Bp_{21}+Ap_{12}$ & 24 & $AB(A+B)-m(A^2+B^2)$\\
8 & $c^3-\tfrac23t_at_b$ & 25 & $Bt_a+At_b-m(t_a+t_b)$\\
9 & $cAB$ & 26 & $m(A-m)(B-m)$\\
10 & $mc(A+B)$ & 27 & $(A+B)(A-m)(B-m)$\\
11 & $c\,t_at_b$ & 28 & $m^2(A+B-2m)$\\
12 & $m\ell$ & 29 & $c\,p_{31}-(Bs_a+As_b)/3$\\
13 & $m(p_{22}-AB/4)$ & 30 & $p_{21}^2+p_{12}^2-\big(Bs_a+As_b-AB(A+B)/4\big)/3$\\
14 & $m\,p_{31}$ & 31 & $p_{21}p_{12}-t_at_b/3$\\
15 & $mc(t_a+t_b)$ & 32 & $(A+B)(p_{22}-AB/4)$\\
16 & $m\ell(A+B)$ & & \\
\bottomrule
\end{tabular}
\caption{The functions in the kernel.}\label{tab:features}
\end{table}

\begin{proposition}\label{prop:expect}
Let $\mu$ be a colour-symmetric fixed point of $\mathcal T_{1/\sqrt2}$ with second moment $m$ and third
moment $T$, and let $a,b$ be independent with law $\mu$. Then
\[
\E\,K(a,b;m)=256\,(x-x_{\mathrm{out}})-\beta m^3-g^*mT,\qquad
g^*=\tfrac{16}3-g_0(2\sqrt2-2)>\tfrac{863}{375000000}.
\]
\end{proposition}

\begin{proof}
At $\lambda=1/\sqrt2$ the fixed-point equations \eqref{eq:moments-out} read $\E F_3=2\sqrt2\,T$ and
$\E F_4=4V$. Hence $\E\phi_0=m(2\sqrt2-2)T$ and $\E\phi_1=2mV-m(4V-2V)=0$. The expectations of
$\phi_2,\dots,\phi_{32}$ vanish: for $\phi_5,\phi_6,\phi_7,\phi_9,\phi_{10},\phi_{11},\phi_{12},\phi_{14},
\phi_{15},\phi_{16},\phi_{22}$ by \eqref{eq:weighted}, for $\phi_{23}$ by the symmetry of the law of $b$,
for $\phi_{24}$ through $\phi_{28}$ by independence and $\E A=\E B=m$, and for the others by
\cref{tab:moments}. Finally $\E[\tfrac23t_at_b]=\tfrac23T^2$ and $\E[\beta mAB]=\beta m^3$; combining
with \eqref{eq:drift} gives the formula. The bound on $g^*$ uses $g_0=32189/5000$ and
$\sqrt2<70711/50000$.
\end{proof}

\begin{proposition}\label{prop:kernel}
Let $\varrho=677/1000$ and $M=8U=1375940705703331/3799912185593856$. Then $K(a,b;m)\ge0$ for all
$a,b\in\varrho(4\Simp-\one)$ and $0\le m\le M$.
\end{proposition}

\begin{proof}
Since $K$ is affine in $m$, it suffices to treat $m=0$ and $m=M$. Every function in $K$ is invariant
under a simultaneous permutation of the colours of $a$ and $b$, so we may assume
$a_0\le a_1\le a_2\le a_3$; the coordinates of $b$ are then sorted by one of the $24$ permutations
$\pi$. Write $a$ and the sorted $b$ in barycentric coordinates $z,w\ge0$ of the ordered simplex with
vertices $\varrho(-1,-1,-1,3)$, $\varrho(-1,-1,1,1)$, $\varrho(-1,\tfrac13,\tfrac13,\tfrac13)$ and $0$, and
split each of the two charts into the cells $z_0\ge z_3$ and $z_3\ge z_0$ (respectively for $w$),
substituting $z_0=z_0'+z_3$ or $z_3=z_3'+z_0$. This gives $24\times4\times2=192$ cases. In each case
\[
G=K\cdot P^5Q^5S^3(UV)^e,\qquad U=\textstyle\sum z_i,\ V=\sum w_i,\ P=3\cdot10^6U,\ Q=3\cdot10^6V,
\]
where $a_i=\alpha_i/P$ and $b_i=\alpha_i'/Q$ with integer linear forms $\alpha_i,\alpha_i'$ in the
cell variables and $S=4PQ+\sum_i\alpha_i\alpha_i'$ is the cleared form of $s$, is a polynomial in the cell variables, and
$S>0$ because $s\ge1-\varrho^2$. Let
\[
\Sigma_1=AB(A+B)-2t_at_b,\ \Sigma_2=(A+B)(AB-c^2),\ \Sigma_3=m(AB-c^2),\ \Sigma_4=AB(A+B)+2t_at_b .
\]
Each $\Sigma_k$ is nonnegative: $\Sigma_2,\Sigma_3$ by Cauchy--Schwarz, and $\Sigma_1,\Sigma_4$ because
$|t_at_b|\le(AB)^{3/2}\le AB(A+B)/2$. For each case there are rationals $\gamma_k\ge0$ such that every
coefficient of
\[
G-\sum_k\gamma_k\,\Sigma_k\,P^5Q^5S^3(UV)^e
\]
is nonnegative; the degree $e$ is $1$ in $176$ cases and $2$ in $16$, $\gamma=0$ in $168$ cases and
$\gamma=(1/3,0,0,0)$ in the remaining $24$. Since the cell variables are nonnegative and the factors
$P,Q,S,U,V$ are positive, $K\ge\sum_k\gamma_k\Sigma_k\ge0$ on every cell.
\end{proof}

\section{Proof of the main theorem}\label{sec:proof}

Let $\lambda=1/\sqrt2$ and let $\mu_\infty$ be the limit law of \cref{lem:fixed}, with moments $m$ and
$T$ and posterior second moment $x$. By \cref{prop:grid,lem:fixed}, $x\le U$, so $m=8x\le M$. By
\cref{lem:support} the law is supported in $\varrho(4\Simp-\one)$, so \cref{prop:kernel} applies to
every pair in its support and $\E K\ge0$. At a fixed point $x_{\mathrm{out}}=x$, so
\cref{prop:expect} gives $\beta m^3+g^*mT\le0$. By \cref{prop:cubic}, $T\ge0$, and therefore $m=0$.
Thus $\mu_\infty$ is the point mass at zero and the process is nonreconstructible at $\lambda=1/\sqrt2$.
For $\lambda<1/\sqrt2$ the same follows from \cref{lem:garble}. Reconstruction for $\lambda>1/\sqrt2$ is
the Kesten--Stigum theorem. \qed

\begin{remark}
Since $\E\phi_{24}=0$ at every colour-symmetric law, the coefficient $g_{24}$ does not affect
\cref{prop:expect}; its value only serves the certificates of \cref{prop:kernel}. The certificates are
not slack: with $\gamma=0$ in the cases where $\gamma\ne0$, or with the shift $1743/480965$ in $g_{24}$ replaced by $-1/10$, some
coefficients become negative.
\end{remark}

\section{The computations}\label{sec:computations}

All computations use exact integer or rational arithmetic; no floating-point value enters a bound.
The grid computation of \cref{prop:grid} handles the $114481$ unordered pairs of the $478$ orbits per
step and stores masses as integers over $2^{48}$. Two independent implementations of
\cref{lem:split} and \eqref{eq:bp} agree on all $500$ intermediate moments and on the final value $U$.
The coefficient comparisons of \cref{prop:cubic} cover $6\times129600$ inequalities. The $192$
certificates of \cref{prop:kernel} were computed with \texttt{python-flint}; each polynomial is rebuilt
from the definitions of \cref{tab:features}, and its value was compared with a direct evaluation of $K$
at rational points. The certificates, the coefficient files and the programs that rebuild and check them
are at \url{https://github.com/yspennstate/potts-four-state-binary-tree}.

\appendix
\section{The coefficients}\label{app:coeff}

\begin{center}\small
\begin{tabular}{rl@{\qquad}rl}
\toprule
$j$ & $g_j$ & $j$ & $g_j$\\
\midrule
0 & $32189/5000$ & 17 & $8086351/894289$\\
1 & $307513/304635$ & 18 & $314989/778090$\\
2 & $1349531/958408$ & 19 & $-293112/998719$\\
3 & $13119169/971828$ & 20 & $390816/998719$\\
4--7 & $0$ & 21 & $2119698/387547$\\
8 & $2046290/998719$ & 22 & $-760908/859123$\\
9 & $0$ & 23 & $-9410/911487$\\
10 & $10803745609/1218540000$ & 24 & $g_{24}$ (below)\\
11 & $649352/585675$ & 25--28 & $0$\\
12 & $-96567/5000$ & 29 & $-1454253/671354$\\
13 & $-33529637/939631$ & 30 & $-1640643/978656$\\
14 & $-96567/5000$ & 31 & $-146556/998719$\\
15 & $1007541/284861$ & 32 & $-505057/629438$\\
16 & $2858842/731359$ & & \\
\bottomrule
\end{tabular}
\end{center}
\begin{align*}
g_{24}&=-\frac{44673597587449418976423934191113}{1462309266444848905334068011200160}+\frac{1743}{480965}\\
&=-\frac{3787526362446847631199691392941357033}{140663915267129350750800004201376990880}\approx-0.026926 .
\end{align*}

\section*{Acknowledgments and funding}
This research was supported by the European Research Council (ERC) under the European Union's Horizon Europe
research and innovation programme (grant agreement No.~101041711), by the Simons Foundation as part of the
Collaboration on the Mathematical and Scientific Foundations of Deep Learning, by Heights Labs, by Convex Nexus
Capital, by the Israel Science Foundation (grant number 2258/19), and by the Israel Science Foundation (ISF Grant
4101/25).

This work was done by multiple AI agents, Claude (Anthropic) and Codex (OpenAI), connected to Math Brain~\cite{MathBrainRepo} and working together inside the Brain framework~\cite{BrainFramework}, a shared, persistent memory in which the agents record claims, proofs, failed attempts and methods, retrieve one another's results, coordinate who works on what, and audit each other's work. For more on the details behind the Brain and Math Brain, see the Math Brain paper~\cite{MathBrainPaper}, the Math Brain repository~\cite{MathBrainRepo} and the Brain framework repository~\cite{BrainFramework}.

\section*{Declaration of competing interest}
The author holds equity in Heights Labs and in Convex Nexus Capital and has a management role at Convex Nexus
Capital. These entities are included in the funding declaration above.

\section*{Code and data availability}
The code and data for this paper are at \url{https://github.com/yspennstate/potts-four-state-binary-tree}. The paper is part of the Math Brain project, whose public repository is \url{https://github.com/yspennstate/math-brain-public}~\cite{MathBrainRepo}.

\section*{Declaration of generative AI and AI-assisted technologies}
The paper comes out of Math Brain, an ongoing project on the use of AI systems for research mathematics. In
October 2026 OpenAI GPT models, running as research tasks of that project, wrote the proofs and the programs that
produce the computer-assisted steps. Claude models (Anthropic), working in Claude Code, audited every piece with
programs of their own, wrote the kernel certificates of \cref{prop:kernel}, assembled the argument, and wrote the
paper. The author directed the work and is responsible for the content of the paper and for decisions concerning
its dissemination.

\begin{thebibliography}{99}
\bibitem{BRZ95} P.~M. Bleher, J.~Ruiz and V.~A. Zagrebnov, On the purity of the limiting Gibbs state for
the Ising model on the Bethe lattice, \emph{J. Stat. Phys.} 79 (1995), 473--482.
\bibitem{BM11} N.~Bhatnagar and E.~Maneva, A computational method for bounding the probability of
reconstruction on trees, \emph{SIAM J. Discrete Math.} 25 (2011), 854--871; arXiv:0903.4812.
\bibitem{EKPS00} W.~Evans, C.~Kenyon, Y.~Peres and L.~J. Schulman, Broadcasting on trees and the Ising
model, \emph{Ann. Appl. Probab.} 10 (2000), 410--433.
\bibitem{KS66} H.~Kesten and B.~P. Stigum, Additional limit theorems for indecomposable multidimensional
Galton--Watson processes, \emph{Ann. Math. Statist.} 37 (1966), 1463--1481.
\bibitem{Mos01} E.~Mossel, Reconstruction on trees: beating the second eigenvalue, \emph{Ann. Appl.
Probab.} 11 (2001), 285--300.
\bibitem{MSS22} E.~Mossel, A.~Sly and Y.~Sohn, Exact phase transitions for stochastic block models and
reconstruction on trees, \emph{Ann. Probab.} 53 (2025), 967--1018; arXiv:2212.03362.
\bibitem{OAI4} OpenAI, The reconstruction threshold for the ferromagnetic four-state Potts model,
preprint, 5 October 2026, \url{https://github.com/openai/math}.
\bibitem{RSZ19} F.~Ricci-Tersenghi, G.~Semerjian and L.~Zdeborov\'a, Typology of phase transitions in
Bayesian inference problems, \emph{Phys. Rev. E} 99 (2019), 042109; arXiv:1806.11013.
\bibitem{Shm3} Y.~Shmalo, The Kesten--Stigum bound is sharp for the three-state Potts model on regular
trees, preprint, 2026, \url{https://doi.org/10.5281/zenodo.24064343}.
\bibitem{Sly11} A.~Sly, Reconstruction for the Potts model, \emph{Ann. Probab.} 39 (2011), 1365--1406;
arXiv:0811.1208.
\bibitem{MathBrainRepo} Y.~Shmalo, \emph{Math Brain: public snapshot, code and tests}, 2026, \url{https://github.com/yspennstate/math-brain-public}.

\bibitem{MathBrainPaper} Y.~Shmalo, \emph{Math Brain: shared, auditable memory for multi-agent mathematical research}, October 2026, \url{https://github.com/yspennstate/math-brain-public/tree/main/papers/math-brain-2026-10-24}.

\bibitem{BrainFramework} Y.~Shmalo, \emph{Brain Framework}, public software repository, 2026, \url{https://github.com/yspennstate/brain-framework-public}.

\end{thebibliography}

\end{document}
